\documentclass[11pt,aspectratio=169]{beamer}
\usepackage[utf8]{inputenc}
\usetheme{Madrid}
%%%%%%%%%%%%%%%%%%%%%%%%%%%%%%%%%%%%%%%%%%%%%%%%%%%%%%%%%%%%%%%%%%%%%%%%%%%%%%%%%%%%%%%%%%%%%%%%%%%%%%%%%%%%%%%%%%%%%%%%%%%%%%%%%%%%%%%
\definecolor{blue_text}{RGB}{0,115,197}
\definecolor{blue1}{RGB}{26,129,203}
\definecolor{blue2}{RGB}{102,171,220}
\definecolor{blue3}{RGB}{154,199,232}

\definecolor{green1}{RGB}{0,128,0}
\definecolor{green2}{RGB}{204,229,204}
\definecolor{green3}{RGB}{229,242,229}

\definecolor{lightblue2}{RGB}{204,227,243}
\definecolor{lightblue3}{RGB}{229,241,249}

\definecolor{orange0}{RGB}{255,116,0}
\definecolor{blue0}{RGB}{0,103,169}%vaidya
\definecolor{green0}{RGB}{0,148,104}


\setbeamercolor*{palette tertiary}{bg=blue1,fg=white}
\setbeamercolor*{palette secondary}{bg=blue2,fg=white}
\setbeamercolor*{palette primary}{bg=blue3,fg=black}

\setbeamercolor{titlelike}{bg=white,fg=blue_text}
\setbeamertemplate{title page}[default][rounded=false, shadow=false]

\setbeamercolor{frametitle}{bg=white,fg=blue_text}

\setbeamercolor{item projected}{bg=blue_text}

\setbeamertemplate{blocks}[rounded]%[shadow]
\setbeamercolor{block title example}{bg=green2,fg=green1}
\setbeamercolor{block body example}{bg=green3,fg=black}

\setbeamercolor{block title}{bg=lightblue2,fg=blue_text}
\setbeamercolor{block body}{bg=lightblue3,fg=black}

\setbeamercolor{button}{bg=blue2,fg=white}

\setbeamertemplate{navigation symbols}{}
%%%%%%%%%%%%%%%%%%%%%%%%%%%%%%%%%%%%%%%%%%%%%%%%%%%%%%%%%%%%%%%%%%%%%%%%%%%%%%%%%%%%%%%%%%%%%%%%%%%%%%%%%%%%%%%%%%%%%%%%%%%%%%%%%%%%%%%

\usepackage{eqnarray,amsmath}
\usepackage{amsfonts}
\usepackage{amssymb}
\usepackage{graphicx}
\usepackage{lmodern} 
\usepackage{bm}
\usepackage{epstopdf}
\usepackage{changepage}
\usepackage{array,booktabs}

\usepackage{pgfplots}
\usepackage{tikz}
\usetikzlibrary{patterns}
\usetikzlibrary{plotmarks}
\usepgfplotslibrary{fillbetween}

\DeclareMathOperator*{\argmax}{arg\,max}
\newcommand*\diff{\mathop{}\!\mathrm{d}}
\DeclareMathOperator{\supp}{supp}
\DeclareMathOperator{\prob}{Pr}

\usepackage{transparent}
\newcommand{\semitransp}[2][35]{\color{fg!#1}#2}

\usetikzlibrary{arrows.meta}

%%%%%%%%%%%%%%%%%%%%%%%%%%%%%%%%%%%%%%%%%%%%%%%%%%%%%%%%%%%%%%%%%%%%%%%%%%%%%%%%%%%%%%%%%%%%%%%%%%%%%%%%%%%%%%%%%%%%%%%%%%%%%%%%%%%%%%%


%====================================================
%========== DEFINITION OF AUTHORS ETC...=============
%====================================================

\author[Ash, Lou, Song]{%
    \texorpdfstring{%
    \begin{columns}
        \column{.3333\linewidth}
        \centering
        Elliot Ash \\\vspace{.5em} \small ETH Zurich
        \column{.3333\linewidth}
        \centering
        Yichuan Lou \\\vspace{.5em} \small UTokyo
        \column{.3333\linewidth}
        \centering
        Lena Song \\\vspace{.5em} \small UIUC
    \end{columns}}{Elliot Ash, Yichuan Lou, Lena Song}
}

\title[AI and Communication]{AI and Communication}
%\subtitle{Subtitle}
\date[Conference Name]{September 2025}




%====================================================
%========== BEGINNING OF DOCUMENT ===================
%====================================================
\begin{document}

\begin{frame}[plain]
    \titlepage
\end{frame}




\begin{frame}{Theoretical Framework: Overview}

\begin{itemize}\setlength{\itemsep}{1em}
    \item Framework based on Benabou and Tirole (2006) to model social signaling with image concerns in the presence of AI senders/receivers
    \item Two novel ingredients: outside option and privacy concerns
    \item Human sender makes two sequential decisions: whether to join the game or not; if he joins, how much effort to exert
    \item Human sender makes his decision based on
    \begin{enumerate}
        \item intrinsic motivation
        \item effort cost
        \item image concerns
        \item privacy concerns
    \end{enumerate}
\end{itemize}
    
\end{frame}



\begin{frame}{Composition of Senders}

\begin{itemize}
    \item Three identity groups: humans, Gen-AI bots, and rule-based bots
    \item Denote by $m_h,m_g,m_r$ the number of different groups of senders 
\end{itemize}
    
\end{frame}



\begin{frame}{Timing}

\begin{enumerate}\setlength{\itemsep}{1em}
    \item A sender is matched with a receiver
    \item The sender's identity (human or bot) and type ($v\in [\underline{v},\overline{v}]$ if human) is initially unknown to the receiver
    \item The human sender chooses whether to participate in the signaling game
    \item If the human sender chooses to not participate, he obtains outside option with payoff $u_0$
    \item If the human sender chooses to participate, he further chooses a signaling action $a\in \mathbb{R}$
    \item The belief about the sender's identity and type is updated
    \item If the sender is a human, his payoff is realized
\end{enumerate}
    
\end{frame}


\begin{frame}{Human Sender}
If a human sender participates and then chooses action $a\in \{a_1,a_2\}$, he obtains
\begin{align*}
    U = va - C(a) + \prob(\text{human}|a)\Big[\underbrace{\mu E(v|a)}_{\text{image concerns}} - \underbrace{\gamma}_{{\text{privacy concern}}}\Big]
\end{align*}
\vspace{1em}

\begin{itemize}
    \item $\prob(\text{human}|a)$ is the posterior probability that the sender is human
    \item $E(v|a)$ is the posterior expectation of the human sender's type $v$ 
    \item $\mu$ captures the strength of the human sender's image concerns
    \item $\gamma$ captures the human sender's concern about privacy leakage, i.e., the risk of revealing his identity (human v.s. bot), which is independent of the belief about $v$
\end{itemize}
\end{frame}


\begin{frame}{Bot Senders}
\begin{itemize}\setlength{\itemsep}{1em}
    \item A bot sender, whether Gen-AI or rule-based, has no private types or preferences
    \item Each bot sender follows a fixed action pattern
\end{itemize}
\vspace{1em}

\begin{exampleblock}{Assumption 1}
    A rule-based bot sender always chooses a detectable action $a_r\notin\{a_1,a_2\}$.
\end{exampleblock}

\begin{exampleblock}{Assumption 2}
    A Gen-AI sender always randomizes between $a_1$ and $a_2$: it plays $a_1$ w.p. $\alpha$ and $a_2$ w.p. $(1-\alpha)$.
\end{exampleblock}
\end{frame}

\begin{frame}{Interior Equilibrium with Cutoffs $v_1<v_2$}

We focus on threshold equilibria characterized by two cutoffs $v_1<v_2$. In such an equilibrium, the human sender chooses
\begin{itemize}
    \item \textbf{Out} if $v\in [\underline{v},v_1)$
    \item (\textbf{In}, $a_1$) if $v\in (v_1,v_2)$
    \item (\textbf{In}, $a_2$) if $v\in (v_2,\bar{v}]$
\end{itemize}

\end{frame}

\begin{frame}{Testing Hypotheses}

\begin{exampleblock}{Hypothesis I}
    Under some conditions \textcolor{red}{(high $\gamma$?)}, the cutoff $v_1$ decreases in the share of Gen-AI bots $\frac{m_g}{m_g+m_h}$. In other words, a higher proportion of Gen-AI bots raises overall human participation.
\end{exampleblock}
\vspace{2em}

\begin{exampleblock}{Hypothesis II}
    Under some conditions, the cutoff $v_2$ decreases in the share of Gen-AI bots $\frac{m_g}{m_g+m_h}$. In other words, a higher proportion of Gen-AI bots raises effort level among participating human senders.
\end{exampleblock}
    
\end{frame}


\end{document}